% !TEX program = lualatex
% =============================================================================
%  Guia Rápido de Uso — Hanaro — modelo AX Academy / IFAM
%
%  COMPILAR COM:  lualatex guia-rapido.tex     (duas vezes: sumário e páginas)
%
%  Todo o estilo mora em axacademy-guia.sty. Aqui ficam só os dados do guia e
%  a ordem dos capítulos, que estão em caps/. As telas ficam em img/.
%
%  O conteúdo é o do Hanaro, o sistema de controle de perdas (scrap) da Linha
%  Final, e serve de gabarito para os outros projetos: telas reais do MVP,
%  oito passos, duas partes com um capítulo cada. Ver README.md.
%
%  Sem número de equipe em nenhum lugar do guia: a autoria é a equipe do Polo
%  INOVA. Subtítulo e tagline vazios tiram o texto da capa, da folha de rosto
%  e da contracapa.
% =============================================================================
\documentclass[11pt,a4paper,oneside,openany]{book}
\usepackage{axacademy-guia}

% ---------- DADOS DESTE GUIA ----------
\guiatitulo{Guia Rápido\\de Uso}
\guiasubtitulo{}
\guiasistema{Hanaro}
\guiatagline{}
\guiaequipe{equipe do Polo INOVA}
\guiaano{2026}
\guiaversao{1.0}
\guiacidade{Manaus, Amazonas --- Brasil}
\guiacliente{LG Electronics do Brasil}
\guiareferencia{}

% ---------- CONTATO (folha de rosto) ----------
% Campo vazio some da folha de rosto, com o ícone. O suporte do Hanaro é feito
% pelo próprio sistema (Capítulo 6), por isso não há e-mail.
\guiaemail{}
\guiasite{www.ifam.edu.br}
\guiaendereco{INOVA --- Polo de Inovação do IFAM\\
    Av. Gov. Danilo de Matos Areosa, 1975 --- Distrito Industrial I\\
    Manaus/AM --- CEP 69075-351}

% ---------- AVISO DA CAPA ----------
% Quebras explícitas: a caixa comporta três linhas curtas.
\guiaaviso{LEIA ANTES DO\\PRIMEIRO ACESSO}

% ---------- MARCA DO SISTEMA ----------
% O Hanaro tem o logotipo inteiro em arquivo (art/marca/hanaro-assinatura.pdf),
% então o ícone fica vazio e o nome não é composto em texto.
\guiamarcaicone{}
\guiamarcatexto{hanaro-assinatura}

% ---------- IMAGEM DA CAPA ----------
% Vazio: capa sem imagem. Para pôr uma, nomeie um arquivo de img/.
\guiacapaimagem{}

% ---------- TEXTO DA CONTRACAPA ----------
% Sai logo abaixo das marcas da AX Academy e do IFAM.
\guiasobreaequipe{Guia de uso rápido baseado no estado do projeto em outubro de 2026}

\begin{document}

% ---------- ABERTURA ----------
\guiacapa
\guiafolhaderosto
\guiaavisodeuso

% ---------- HISTÓRICO DE REVISÕES ----------
\begin{guiahistorico}
    \guiarevisao{1.0}{Outubro de 2026}{AX Academy}{Emissão inicial.}
\end{guiahistorico}

\begin{remark}[info]
Este guia cobre o uso básico do Hanaro pelo perfil Gestor. Quaisquer outros dados e
informações adicionais do projeto, por favor consulte o manual completo.
\end{remark}

\clearpage

\guiasobreax

\guiapagina{Sobre o sistema}{%
    O Hanaro é o sistema web de controle das perdas produtivas (\emph{scrap}) da
    Linha Final. Ele recebe o relatório \emph{Other Account Transaction Text}
    exportado do GERP, consolida a base de scrap e publica o IF Cost no Dashboard e
    no Modo TV, exibido nas telas da fábrica.

    Antes dele, o analista de Produção reprocessava à mão a planilha enviada pela
    Qualidade, de 1,5 a 2 horas por dia, para converter os valores, recortar o
    escopo e eleger os ofensores do período.

    O fluxo básico tem oito passos: entrar no sistema, enviar o relatório do GERP,
    conferir a execução, ler o Dashboard, revisar a Base de Scrap, consultar a
    análise de uma ocorrência, aplicar um modelo de análise em lote e exibir o Modo
    TV. A \autoref{sec:fluxo} apresenta esse fluxo em BPMN; o \autoref{cap:passos}
    descreve cada passo com a tela correspondente.
}

\guiasumario

% ---------- PARTE I ----------
% Um capítulo só por parte; dentro dele, só seções.
\part{Visão Geral}
% =============================================================================
%  Capítulo 1 — Antes de Começar (Parte I, Visão Geral)
%  Capítulo único da parte: público, requisitos de acesso, como ler o guia e
%  o fluxo de uso em BPMN, com a legenda e o mapa passo -> tela -> página.
% =============================================================================
\chapter{Antes de Começar}

\section{Público e escopo}

Este guia se destina ao usuário com perfil \textbf{Gestor}, que mantém a base de
scrap atualizada e acompanha o IF Cost da Linha Final. A leitura completa leva cerca
de quinze minutos.

\begin{itemize}
    \item \textbf{Coberto:} acesso, envio do relatório do GERP, conferência da
          execução, Dashboard, Base de Scrap, análise de uma ocorrência, aplicação de
          modelo em lote e Modo TV.
    \item \textbf{Não coberto:} relatórios de análise, alertas e regras de
          notificação, configurações (classificações de material, tipos de scrap,
          programa do ano) e administração de usuários.
\end{itemize}

\section{Requisitos de acesso}

\begingroup
\guiazebra
\small
\noindent\begin{longtable}{@{}L{4.0cm}L{11.5cm}@{}}
    \toprule
    \guiacabtabela
    \thead{Item} & \thead{Requisito} \\
    \midrule
    \endhead
    \bottomrule
    \endlastfoot
    Navegador           & Navegador atualizado. As telas deste guia foram capturadas no
                          Google Chrome \\
    Endereço do sistema & Informado pela equipe na liberação do acesso \\
    Credenciais         & Usuário e senha do Hanaro, cadastrados pelo administrador \\
    Perfil              & Gestor \\
    Arquivo de entrada  & Relatório \emph{Other Account Transaction Text} exportado do
                          GERP, em texto separado por tabulação (TSV), com codificação
                          UTF-8 ou CP1252 \\
    Modo TV             & TV ou monitor da linha com navegador, para a exibição em tela
                          cheia \\
\end{longtable}
\endgroup

\begin{remark}
O usuário só entra no Hanaro depois de cadastrado e ativado pelo administrador do
sistema. Sem esse cadastro, a tela de acesso recusa as credenciais.
\end{remark}

\section{Como ler este guia}

O \autoref{cap:passos} descreve o uso em \textbf{passos}. Cada passo ocupa uma linha
da página, dividida em duas colunas:

\begin{itemize}
    \item \textbf{à esquerda, a tela} do sistema, exatamente como o usuário a vê;
    \item \textbf{à direita, o texto} do passo: o que fazer, em que ordem e o que o
          sistema responde.
\end{itemize}

Os números em vermelho sobre a tela, como \marca{1} e \marca{2}, indicam os elementos
citados no texto ao lado, na mesma ordem. O número do passo, no selo
\colorbox{guiavermelho}{\color{white}\guiahead\bfseries\fontsize{7.5}{7.5}\selectfont PASSO 1},
é o mesmo que aparece no diagrama BPMN da \autoref{sec:fluxo}.

Botões e opções da interface aparecem em destaque, como \botao{Entrar}. Um caminho de
menu é escrito como \caminho{Menu\seta Execuções}.

As caixas ao longo do texto têm três tipos:

\begin{remark}
\textbf{Atenção.} Condição que impede o passo de ser concluído ou que exige
correção.
\end{remark}

\begin{remark}[dica]
\textbf{Dica.} Atalho ou forma mais rápida de concluir o passo.
\end{remark}

\begin{remark}[info]
\textbf{Informação.} Comportamento do sistema que convém conhecer, sem ação
necessária.
\end{remark}

% ---------------------------------------------------------------------------
%  Fluxo de uso em BPMN
%
%  O diagrama usa os estilos `bpmn ...' do axacademy-guia.sty. As raias são
%  COLUNAS e o fluxo desce: assim o diagrama cabe em pé, nos 16 cm da largura
%  do texto. Coordenadas em cm, com y = 0 no topo das raias.
%
%  A seção abre página nova: o diagrama vai em `figure[H]', que não flutua, e
%  precisa da página inteira. Com oito passos, as tarefas têm uma linha só e
%  as linhas do diagrama ficam a 1,15 cm uma da outra; com mais passos,
%  divida o fluxo em dois diagramas.
% ---------------------------------------------------------------------------
\clearpage
\section{Fluxo de uso}\label{sec:fluxo}

O fluxo básico vai do acesso ao Modo TV. A raia \textbf{Gestor} reúne o que o
usuário faz; a raia \textbf{Hanaro}, o que o sistema faz em resposta. As tarefas com
borda vermelha são os passos deste guia: o número no canto é o do passo e leva à
página em que ele é descrito.

\begin{figure}[H]
\centering
\begin{tikzpicture}[x=1cm,y=1cm,
        bpmn tarefa/.append style={text width=4.4cm,minimum height=0.72cm}]
    % A caixa de medida é a das raias, exata: sem ela, a espessura do contorno
    % passa 0,7 pt da largura do texto.
    \useasboundingbox (0,0) rectangle (16,-13.35);
    % ---- raias ----
    \bpmnraia{0}{8}{13.35}{Gestor}
    \bpmnraia{8}{8}{13.35}{Hanaro}

    % ---- gestor ----
    \node[bpmn inicio] (ini) at (4,-1.05) {};
    \node[bpmn rotulo,anchor=west] at ([xshift=4pt]ini.east)
        {Relatório do GERP\\disponível};
    \bpmnpasso{p1}{4,-2.15}{1}{Entrar com usuário e senha}
    \bpmnpasso{p2}{4,-4.45}{2}{Enviar o relatório do GERP}
    \bpmnpasso{p3}{4,-6.75}{3}{Conferir a execução}
    \bpmnpasso{p4}{4,-7.90}{4}{Ler o Dashboard}
    \bpmnpasso{p5}{4,-9.05}{5}{Revisar a Base de Scrap}
    \bpmnpasso{p6}{4,-10.20}{6}{Consultar a análise do caso}
    \bpmnpasso{p7}{4,-11.35}{7}{Aplicar o modelo em lote}
    \bpmnpasso{p8}{4,-12.50}{8}{Exibir o Modo TV}

    % ---- sistema ----
    \node[bpmn tarefa] (s1) at (12,-2.15) {Validar as credenciais};
    \bpmngateway{g1}{12,-3.30}
    \node[bpmn rotulo,anchor=west] at ([xshift=3pt]g1.east) {Credenciais\\válidas?};
    \node[bpmn tarefa] (s2) at (12,-4.45) {Abrir o Dashboard};
    \node[bpmn tarefa] (s3) at (12,-5.60) {Processar o arquivo enviado};
    \node[bpmn tarefa] (s4) at (12,-6.75) {Publicar o snapshot da base};
    \node[bpmn tarefa] (s5) at (11.6,-12.50) {Girar as telas do Modo TV};
    \node[bpmn fim] (fim) at (14.95,-12.50) {};
    \node[bpmn rotulo,anchor=south,align=center] at ([yshift=3pt]fim.north)
        {Indicadores\\na fábrica};

    % ---- fluxos ----
    \draw[bpmn fluxo] (ini) -- (p1);
    \draw[bpmn fluxo] (p1) -- (s1);
    \draw[bpmn fluxo] (s1) -- (g1);
    \draw[bpmn fluxo] (g1) -- node[bpmn rotulo,anchor=west,xshift=2pt] {Sim} (s2);
    \draw[bpmn fluxo] (g1.west) -- node[bpmn rotulo,anchor=south,pos=0.25] {Não}
        (4,-3.30) -- (p1.south);
    \draw[bpmn fluxo] (s2) -- (p2);
    \draw[bpmn fluxo] (p2.south) |- (s3.west);
    \draw[bpmn fluxo] (s3) -- (s4);
    \draw[bpmn fluxo] (s4) -- (p3);
    \draw[bpmn fluxo] (p3) -- (p4);
    \draw[bpmn fluxo] (p4) -- (p5);
    \draw[bpmn fluxo] (p5) -- (p6);
    \draw[bpmn fluxo] (p6) -- (p7);
    \draw[bpmn fluxo] (p7) -- (p8);
    \draw[bpmn fluxo] (p8) -- (s5);
    \draw[bpmn fluxo] (s5) -- (fim);
\end{tikzpicture}
\caption{Fluxo básico de uso do Hanaro pelo Gestor, em BPMN. Os números indicam os
passos do \autoref{cap:passos}.}
\label{fig:bpmn}
\end{figure}

\section{Leitura do diagrama}

\begingroup
\small
\renewcommand{\arraystretch}{1.6}
\noindent\begin{tabularx}{\textwidth}{@{}M{2.4cm}X@{}}
    \toprule
    \guiacabtabela
    \thead{Símbolo} & \thead{Significado} \\
    \midrule
    \tikz[baseline=-0.6ex]\node[bpmn inicio,minimum size=0.42cm]{};
        & Início do fluxo: o evento que leva o usuário ao sistema. \\
    \tikz[baseline=-0.6ex]\node[bpmn tarefa,text width=1.2cm,minimum height=0.5cm,
            font=\fontsize{7}{8}\selectfont]{Tarefa};
        & Tarefa executada pelo usuário ou pelo sistema, conforme a raia em que está. \\
    \tikz[baseline=-0.6ex]{\node[bpmn tarefa,draw=guiavermelho,line width=1.1pt,
            text width=1.2cm,minimum height=0.5cm,font=\fontsize{7}{8}\selectfont] (t) {Passo};
            \node[guia marcador,minimum size=11pt,font=\guiahead\bfseries\fontsize{7}{7}\selectfont]
            at (t.north west) {3};}
        & Tarefa descrita neste guia. O número é o do passo; no PDF, ele leva à página
          correspondente. \\
    \tikz[baseline=-0.6ex]{\bpmngateway{g}{0,0}}
        & Decisão. O fluxo segue por um único caminho, conforme a condição escrita ao
          lado. \\
    \tikz[baseline=-0.6ex]\draw[bpmn fluxo] (0,0) -- (1.2,0);
        & Sequência: a próxima tarefa só começa quando a anterior termina. \\
    \tikz[baseline=-0.6ex]\node[bpmn fim,minimum size=0.42cm]{};
        & Fim do fluxo. \\
    \bottomrule
\end{tabularx}
\endgroup

\section{Mapa dos passos}

\begingroup
\guiazebra
\small
\noindent\begin{longtable}{@{}P{1.3cm}L{3.2cm}L{8.4cm}P{1.4cm}@{}}
    \toprule
    \guiacabtabela
    \thead{Passo} & \thead{Tela} & \thead{Ação principal} & \thead{Página} \\
    \midrule
    \endhead
    \bottomrule
    \endlastfoot
    \marca{1} & Acesso           & Informar usuário e senha e clicar em \botao{Entrar}       & \pagina{passo:1} \\
    \marca{2} & Execuções        & Selecionar o arquivo do GERP e clicar em
                                   \botao{Processar relatório}                               & \pagina{passo:2} \\
    \marca{3} & Execuções        & Conferir o status e o snapshot da execução                & \pagina{passo:3} \\
    \marca{4} & Dashboard        & Ler o IF Cost acumulado, o target e os rankings           & \pagina{passo:4} \\
    \marca{5} & Base de Scrap    & Filtrar e revisar as ocorrências de scrap                 & \pagina{passo:5} \\
    \marca{6} & Análise de scrap & Ler o custo, a contenção e as evidências da ocorrência    & \pagina{passo:6} \\
    \marca{7} & Base de Scrap    & Marcar as ocorrências e aplicar o modelo de análise       & \pagina{passo:7} \\
    \marca{8} & Modo TV          & Abrir o Modo TV e controlar a rotação das telas           & \pagina{passo:8} \\
\end{longtable}
\endgroup

\fimdecapitulo


% ---------- PARTE II ----------
\part{Passo a Passo}
% =============================================================================
%  Capítulo 2 — Uso Básico do Hanaro (Parte II, Passo a Passo)
%  Capítulo único da parte, só em seções: os oito passos, as colunas do
%  histórico de execuções, as telas do Modo TV e o encerramento da sessão.
%
%  Cada \passo é uma linha de duas colunas: a tela à esquerda e o texto à
%  direita, e conta como uma seção do capítulo (2.1, 2.2...). O argumento
%  opcional lista as posições dos marcadores sobre a tela, como x/y em fração
%  da imagem, com origem no canto SUPERIOR esquerdo. A numeração segue a
%  ordem da lista e casa com a lista `elementos'.
%
%      \passo[0.335/0.525, 0.665/0.649]{01-login}{Título}{Texto ...}
% =============================================================================
\chapter{Uso Básico do Hanaro}\label{cap:passos}

% ---------- acesso e atualização da base ----------
\passo[0.335/0.525, 0.665/0.649, 0.335/0.772]{01-login}{Entrar no Hanaro}{%
    Abra o endereço do Hanaro no navegador. A tela de acesso é a primeira a aparecer.
    \begin{elementos}
        \item \campo{Usuário}: o nome de usuário cadastrado pelo administrador. Não é
              o e-mail.
        \item \campo{Senha}: a senha do usuário. O ícone de olho, no fim do campo,
              mostra o que foi digitado.
        \item \botao{Entrar}: valida as credenciais e abre o Dashboard.
    \end{elementos}
    Credenciais inválidas mantêm o usuário nesta tela.
}

\passo[0.115/0.181, 0.845/0.146, 0.31/0.521, 0.69/0.669]{02-enviar-relatorio}{Enviar o relatório do GERP}{%
    A base de scrap é atualizada pelo relatório \emph{Other Account Transaction
    Text}. Exporte-o do GERP antes deste passo.
    \begin{elementos}
        \item \caminho{Execuções}, no menu lateral: abre a tela
              \textbf{Execuções e atualização}.
        \item \botao{Enviar relatório}: abre a janela \textbf{Atualizar base por
              relatório}.
        \item \textbf{Área do arquivo}: um clique abre a escolha do arquivo exportado
              do GERP. O nome do arquivo escolhido aparece no lugar de
              \emph{Selecionar relatório GERP}.
        \item \botao{Processar relatório}: envia o arquivo. O Hanaro registra uma
              nova execução, conferida no passo seguinte.
    \end{elementos}
    \botao{Cancelar} fecha a janela sem enviar.
}

\begin{remark}
O arquivo precisa ser o \emph{Other Account Transaction Text}, em texto separado por
tabulação (TSV), com codificação UTF-8 ou CP1252, como indica a própria janela.
Planilha salva em outro formato não serve.
\end{remark}

\passo[0.148/0.378, 0.98/0.378, 0.871/0.242]{03-execucoes}{Conferir a execução}{%
    A execução nova entra no topo do \textbf{Histórico de rotinas}.
    \begin{elementos}
        \item \textbf{Linha da execução}: o identificador (\texttt{EXE-...}), a
              origem e o gatilho. Um clique na linha abre os \textbf{Detalhes da
              execução}, com o arquivo de origem e as contagens.
        \item \textbf{Status}: \emph{Na fila} enquanto a execução aguarda o
              processamento e \emph{Concluído} ao final. Ao lado, o
              \textbf{Snapshot} passa de \emph{Não publicado} a \emph{Publicado}.
        \item \botao{Filtrar}: restringe o histórico.
    \end{elementos}
    As colunas estão descritas na \autoref{tab:historico}.
}

\begin{remark}[info]
O Dashboard passa a refletir o arquivo novo quando o snapshot da execução fica
\emph{Publicado}. Até lá, ele mostra a base da execução anterior.
\end{remark}

\section{Colunas do histórico}

% Tabela curta: vai inteira, em `table', para não partir entre páginas.
\begin{table}[H]
\caption{Colunas do histórico de rotinas.}\label{tab:historico}
\guiazebra
\small
\noindent\begin{tabularx}{\textwidth}{@{}L{4.4cm}X@{}}
    \toprule
    \guiacabtabela
    \thead{Coluna} & \thead{Conteúdo} \\
    \midrule
    \textbf{Execution ID / Processo} & Identificador da execução e nome do relatório
                                       processado \\
    \textbf{Origem}                  & Sistema de onde veio o arquivo: GERP \\
    \textbf{Gatilho}                 & Como a execução começou. \emph{Manual}: pelo
                                       envio do relatório no passo~2 \\
    \textbf{Início}, \textbf{Fim} e \textbf{Duração}
                                     & Data e hora do processamento e o tempo que ele
                                       levou \\
    \textbf{Recebidos}               & Registros lidos no arquivo \\
    \textbf{Válidos}                 & Registros aceitos na base de scrap \\
    \textbf{Rejeitados}              & Registros recusados na validação \\
    \textbf{Snapshot}                & \emph{Publicado} quando a base processada já
                                       alimenta os indicadores; \emph{Não publicado}
                                       antes disso \\
    \textbf{Status}                  & \emph{Na fila} enquanto aguarda;
                                       \emph{Concluído} ao final \\
    \bottomrule
\end{tabularx}
\end{table}

\begin{remark}
Com \textbf{Rejeitados} acima de zero, parte do arquivo ficou fora da base. Abra os
detalhes da execução e acione o suporte em \caminho{Ajuda e suporte}, no pé do menu
lateral, com o identificador da execução.
\end{remark}

% ---------- indicadores e ocorrências ----------
\passo[0.30/0.114, 0.36/0.214, 0.66/0.114, 0.45/0.322, 0.924/0.503, 0.27/0.725]{04-dashboard}{Ler o Dashboard}{%
    O Dashboard é a tela que abre depois do acesso. No alto, à direita,
    \emph{Atualizado às} indica a hora da última atualização.
    \begin{elementos}
        \item \textbf{Impacto absoluto}: chave que alterna os indicadores entre o IF
              Cost em dólar e o Scrap Rate, em percentual.
        \item \textbf{Indicadores}: IF Cost acumulado, target, atingimento do target,
              diferença para a meta e variação sobre o mesmo acumulado do ano
              anterior.
        \item \textbf{Controles}: o botão com o ícone de tela abre o Modo TV
              (passo~8); \botao{YTD anterior} escolhe a base de comparação;
              \botao{Filtros} recorta os dados; \textbf{Valor monetário} alterna
              entre dólar e quantidade, em unidades.
        \item \textbf{Target × realizado mensal}: realizado do ano, referência do ano
              anterior e target, mês a mês.
        \item \textbf{Top 10 produtos / áreas}: o IF Cost por produto ou área.
        \item \textbf{Top 10 linhas} e \textbf{Top 10 modelos}: os rankings de
              ofensores. Mais abaixo, componentes e ofensores.
    \end{elementos}
}

\begin{remark}[dica]
Cada quadro tem o seu botão \botao{Filtros} e um botão de tabela, que abre os mesmos
números em linhas e colunas. Passar o cursor sobre um ponto do gráfico mostra os
valores do mês.
\end{remark}

\passo[0.12/0.178, 0.555/0.194, 0.135/0.337, 0.70/0.1]{05-base-scrap}{Revisar a Base de Scrap}{%
    A lista traz as ocorrências consolidadas pelas execuções. Ao lado do título, o
    total de ocorrências, o valor em dólar e a quantidade.
    \begin{elementos}
        \item \caminho{Base de Scrap}, no menu lateral: abre a tela.
        \item \textbf{Ordem e filtros}: \botao{Data da transação} escolhe o campo da
              ordem e \botao{Decrescente} a inverte; a lista \botao{Todos} filtra
              pelo status da análise; \botao{Filtrar} abre os demais filtros.
        \item \textbf{Ocorrência}: organização e produto, data, item, descrição,
              ordem, quantidade, valor, tipo de scrap, status da análise
              (\emph{Não analisado} ou \emph{Revisado}), responsável e relatório. Um
              clique abre a análise (passo~6).
        \item \botao{Aplicar Modelo} e \botao{Selecionar itens}: levam uma análise já
              feita a várias ocorrências (passo~7).
    \end{elementos}
}

\passo[0.39/0.138, 0.195/0.349, 0.195/0.507, 0.195/0.652, 0.577/0.87]{06-ocorrencia}{Consultar a análise de uma ocorrência}{%
    A janela \textbf{Análise de scrap} reúne o que se sabe da ocorrência. Role a
    janela para ver todas as seções.
    \begin{elementos}
        \item \textbf{Status da análise}: \emph{Revisado}, como nesta ocorrência, ou
              \emph{Não analisado}.
        \item \textbf{Custo}: tipo de custo, modelo e quantidade. Acima, \textbf{Onde
              e quando} traz a data e a hora.
        \item \textbf{Contenção imediata}: a ação tomada na linha para conter o
              problema.
        \item \textbf{Evidências e fotos}: as imagens anexadas à análise.
        \item \botao{Criar relatório desta análise}: abre um relatório de análise a
              partir desta ocorrência.
    \end{elementos}
    O coração, no alto, guarda a análise como modelo para o passo seguinte.
}

\passo[0.565/0.09, 0.145/0.494, 0.313/0.632, 0.685/0.768]{07-modelo-lote}{Aplicar um modelo em lote}{%
    Quando o mesmo problema se repete em várias ocorrências, a análise de uma serve
    de modelo para as outras.
    \begin{elementos}
        \item \textbf{Modelo em uso}: em \botao{Aplicar Modelo}, escolha o modelo em
              \textbf{Modelos salvos} e clique em \botao{Usar este modelo}. A faixa
              no alto confirma a escolha.
        \item \textbf{Ocorrências marcadas}: com \botao{Selecionar itens}, marque as
              ocorrências e clique em \botao{Aplicar modelo aos itens selecionados},
              acima da lista. Abre esta janela.
        \item \textbf{Reaproveitar imagens do relatório de referência}: copia as
              fotos do modelo para todas as ocorrências marcadas.
        \item \botao{Aplicar modelo aos itens selecionados}: grava a análise do
              modelo nas ocorrências marcadas.
    \end{elementos}
}

\begin{remark}
Ocorrências inativas ou que já têm relatório são ignoradas na aplicação, como avisa a
própria janela. Copiar as fotos replica as mesmas evidências em todas as ocorrências
marcadas: use só quando as fotos valem para todas.
\end{remark}

% ---------- Modo TV ----------
\passo[0.395/0.069, 0.712/0.083, 0.2/0.865, 0.5/0.865, 0.785/0.865]{08-modo-tv}{Exibir o Modo TV}{%
    O Modo TV mostra os indicadores em tela cheia, para as TVs da fábrica. Ele abre
    pelo botão com o ícone de tela, nos controles do Dashboard (passo~4).
    \begin{elementos}
        \item \textbf{Tela em exibição}: o título da tela. Na tela Fábrica,
              \botao{Mês} e \botao{Ano} escolhem o período.
        \item \textbf{Dados até}: a data do último dado da base e, ao lado, a data e
              a hora atuais.
        \item \textbf{Telas}: \botao{Executivo}, \botao{Fábrica}, \botao{Ofensores} e
              \botao{Prioridades}, descritas na \autoref{tab:telastv}. Um clique leva
              direto à tela.
        \item \textbf{Rotação}: as setas passam de tela, \botao{Pausar} para a
              rotação automática e \textbf{Intervalo} define o tempo de cada tela.
        \item \botao{Tema Claro}, \botao{Tela cheia} e \botao{Sair}: trocam o fundo
              escuro pelo claro, ocupam a tela inteira e voltam ao Dashboard.
    \end{elementos}
}

\begin{remark}[dica]
Na TV da linha, deixe a \textbf{Rotação automática} ligada e use \botao{Tela cheia}:
as quatro telas se alternam sem intervenção. O contador no fim da barra, como
\emph{2/4}, indica a tela atual.
\end{remark}

\section{Telas do Modo TV}

% Tabela curta: vai inteira, em `table', para não partir entre páginas.
\begin{table}[H]
\caption{As quatro telas do Modo TV.}\label{tab:telastv}
\guiazebra
\small
\noindent\begin{tabularx}{\textwidth}{@{}L{2.4cm}L{3.9cm}X@{}}
    \toprule
    \guiacabtabela
    \thead{Botão} & \thead{Título da tela} & \thead{O que mostra} \\
    \midrule
    \textbf{Executivo}   & Visão Executiva
                         & IF Cost acumulado, mesmo período do ano anterior, variação e
                           target acumulado; evolução do IF Cost no ano, no ano anterior
                           e no target; os três componentes mais afetados \\
    \textbf{Fábrica}     & Ranking de ocorrências por linha
                         & Ocorrências do período, linha com mais ocorrências, linhas
                           monitoradas e período anterior; ranking das linhas de
                           montagem, com a variação; setores com mais ocorrências \\
    \textbf{Ofensores}   & Principais ofensores do período
                         & Os cinco defeitos e as oito linhas de maior IF Cost; o part
                           number e o defeito mais críticos \\
    \textbf{Prioridades} & Ocorrências prioritárias
                         & As ocorrências críticas de maior custo, com a descrição;
                           situação das revisões (pendentes, em revisão e
                           justificadas) e categorias predominantes \\
    \bottomrule
\end{tabularx}
\end{table}

\section{Encerramento da sessão}

O nome do usuário fica no pé do menu lateral. O botão \botao{\faEllipsisV}, ao lado
dele, leva ao \textbf{Perfil}. Ali, a seção \textbf{Sessões ativas} lista os
dispositivos conectados à conta: \botao{Encerrar sessão} desconecta um deles e
\botao{Encerrar outras sessões} desconecta todos, menos o atual.

Em computador compartilhado, encerre a sessão ao terminar e feche o navegador.

\fimdecapitulo


% ---------- FECHAMENTO ----------
\guiacontracapa

\end{document}
